\chapter*{如何使用本书，以及旧稿是否完善}
\addcontentsline{toc}{chapter}{如何使用本书，以及旧稿是否完善}
本书整理此前调研的 S1--S17 与主动注入补充文献 P1--P5，共22项。编号沿用原调研，章节按照学习依赖重新排列。S1是教程，S14--S17是统计推断方法或教程，S10--S11是运行域应用；因此，22项文献不能表述为22种可直接比较的拓扑恢复算法。

对上一版《配电网拓扑辨识文献逐篇解析》的审查结论是：\textbf{它适合研究选题和数学边界审阅，但作为算法教材还不完善。}旧稿已经解释共同路径模型、隐藏节点、主动探测和统计保证之间的区别，也标出了全文访问限制。然而，读者仍需自行补全若干关键连接：一个量测文件怎样变成设计矩阵，目标函数怎样变成逐步更新，分组后哪些矩阵和节点集合发生变化，以及什么情况下应当停止并返回“不确定”。本版围绕这些连接重写，保留上一版以便对照。

\begin{longtable}{p{0.17\textwidth}p{0.34\textwidth}p{0.39\textwidth}}
\toprule 审查维度 & 旧稿的实际缺口 & 本版的处理方式\\\midrule
\endhead
算法入口 & 不同论文的量测、符号及未知量分散在段落中 & 在各方法中明确输入、输出、参数和适用数据；先统一物理符号，再说明论文局部符号。\\
步骤闭合 & 若干条目止于目标函数或概念路线 & 补充状态变量、逐步更新、候选管理、停止条件；缺原文细节时显式标注。\\
推导连续性 & 有边界反例，但部分核心更新未从目标导出 & 补充最小二乘、树距离、EM、ADMM、MH、分位数及鲁棒几何的中间推导。\\
算例可操作性 & 数值例子较少，读者不易手算一轮 & 增加贯穿三叶树，以及与各算法匹配的小规模数值轨迹。\\
结论归属 & 独立审查与原文结论在连续段落中交替 & 明确区分原文流程、等价代数整理和教学变体；把访问与公式审计集中到附录。\\
复现证据 & 没有独立重跑论文实验 & 本版验证教学公式与排版，仍不声称复现实验；全文缺口继续保留。\\
\bottomrule
\end{longtable}

\section*{三种算法说明必须分别理解}
\textbf{原文流程重述}意味着已有正文足以核对该步骤的来源，仍可能统一变量名、调整叙述顺序。\textbf{等价代数整理}意味着从已核对的模型或目标推导出便于实现的形式；它并不等于作者公开代码的逐行复现。\textbf{教学变体或背景模块}意味着为解释共同原理而明确选择了模型、求解器或停止准则；它可以独立实现，但不能用对应论文名称冒充论文基线。

本版尤其不能填平 S2、S4、S6 的正文证据缺口。对 S2 和 S4，能做的是身份核实、已确认路线、公共机制教学和待核对清单。S6无法仅由题名补出神经网络层数、训练损失或数据划分。其余条目也必须按各节的来源说明理解，不能因为出现了伪代码就认为已经独立复现。

\section*{建议的阅读路径}
首次学习可先读第1章，亲手验证共同路径矩阵和距离；再读第2章 S3 的矩阵估计、分组与回溯；随后按数据条件进入数据质量、主动探测和统计推断章节。若研究目标是末端 P/Q/V 与隐藏零注入节点，首先掌握“回归可辨性”和“树表示可辨性”这两个不同问题。若研究目标是可信决策，再学习统计覆盖与运行域几何，避免把预测误差小直接解释为拓扑正确或运行安全。

每节的数值例均为教学构造，除非明确写作原文实验。公式中的充分条件只在列出的模型假设下成立。本版不新增原论文实验成绩，不把经验优势改写成普遍定理，也不把有限候选最优写成全结构域最优。
